How do simple local interactions give rise to collective behavior that no single constituent possesses? This question, which I first met as \textit{emergent phenomena} (``More is Different''), is what draws me to graduate study. I want to spend my PhD training in condensed matter theory, focusing on strongly correlated systems and quantum spin liquids, and to become an independent researcher who can find and solve problems in this field.

I first saw how deep the question runs when I attended Prof.~Yi-Ping Huang's talk on quantum spin liquids, phases with no local order parameter whose properties are encoded in long-range entanglement. I joined his group as an undergraduate. My first project was to express the ground state of the Kitaev honeycomb model as a tensor network, and I taught myself frustrated magnetism and tensor-network methods along the way. The research topic was not assigned to me. It took shape through repeated discussions with my advisor, and that taught me to formulate my own questions and defend them. The project earned me an undergraduate research scholarship from Taiwan's National Science and Technology Council (NSTC).

I stayed in the group for my master's degree and began collaborating with Dr.~Sungkit Yip at Academia Sinica, whom Prof.~Huang introduced to me during a 2024 summer internship. We began with the entanglement structure of the Kitaev ground state, working up from a two-spin model to a one-dimensional chain and then to the full honeycomb lattice. Dr.~Yip taught me to \textit{always start from a simple example}, and I have used that principle in every project since.

This work grew into a paper on which I am first author, published in \textit{Physical Review B}. The known entanglement properties of the Kitaev model are computed in an enlarged Majorana Hilbert space with gauge redundancy, so their physical origin is hidden. We worked directly in the physical spin Hilbert space and derived the \textit{exact} density matrix of the Kitaev spin liquid in terms of spin-string operators. It is organized by equivalence classes of strings that depend only on their endpoints. Combining this structure with the Gauss law of the emergent $\mathbb{Z}_2$ gauge theory and the model's exact 1-form Wilson symmetry, we showed that the reduced density matrix splits into $2^{|\partial A|/2-1}$ identical blocks. This solves the entanglement spectrum exactly. The blocks account for its extensive degeneracy and for the gauge part of the entanglement entropy. Our formalism also works for subsystems with an odd number of sites, which is a breakthrough since a Majorana system is ill-defined in such a case. It also reveals a link between the entanglement spectrum and the subsystem fermion parity. {\color{red} my own contribution, derivations, proofs, or numerical checks carried out.} This project covered lattice gauge theory, theory of higher-form symmetry, and quantum-information theory, and it became a valuable experience where I applied the knowledge I acquired in courses to an open problem. My numerical training in symmetry-resolved exact diagonalization and the Lanczos method helped me verify my analytical results.

Since then, I have extended this line of work to quantum magic. I am studying the stabilizer R\'enyi entropy of the Kitaev model, starting with the one-dimensional chain and building a numerical library for it. The density-matrix and string-operator framework looks like a natural way to approach the magic of a reduced density matrix. This is the kind of question that connects condensed matter physics and quantum information, which is where I want to work.

At UCSB, I want to move beyond the Kitaev model to a wider range of strongly correlated systems. My research interests align with the work of Prof.~Leon Balents on frustrated magnets and spin liquids, Prof.~Cenke Xu on topological phases and gauge theories, and Prof.~Chetan Nayak on topological order and quantum information. {\color{red}one sentence per professor naming a specific paper or question that connects to your work?} The Kavli Institute for Theoretical Physics, with its programs and visitors, would also give me a chance to work among leading physicists. My long-term goal is a career in academic research, where I can continue to ask fundamental questions about quantum matter and train the next generation of students.

%%%%%%%%%%%%%%%%%%%%%%%%%%%%%%%%%%%%%%%%%%%%%%%

% Becoming a distinguished physicist and exploring fascinating physical phenomena have always been my goals, and this is exactly why I want to pursue higher education, because having a PhD in physics is necessary for this journey. However, merely a doctorate is far from enough. I ought to earn a degree at a reputable university with a strong condensed matter/quantum matter theory research community, which aligns with my current research interests. Among the universities worldwide, I believe UCSB is the best place to train my research skills, and this is why I’m applying for a PhD in physics, hoping to become a graduate student at UCSB.

% When it comes to my favorite subject, I will always say condensed matter theory, especially strongly correlated systems and frustrated magnets. I found this field fascinating when I attended Prof.~Yi-Ping Huang’s talk about his early work on quantum spin liquids. It was also the first time I learned about \textit{emergent phenomena}—a type of phenomenon caused by collective interactions among a large number of individuals, resulting in an ordered pattern on a larger scale (or, to be short, More is Different). To explore more in this field, I joined Prof.~Huang’s research group to do undergraduate research and began studying quantum spin liquids. 

% I learned a lot of professional knowledge and how a physicist works from my advisor during my undergraduate research. My first project was studying the ground state of the Kitaev honeycomb model and trying to figure out a way to express it using a tensor network, a mathematical formulation that can not only improve numerical efficiency but also reflect the rich entanglement structure of a quantum state. From this project, I learned about frustrated magnets and tensor networks through self-study. More importantly, my advisor didn’t directly give me this research topic. Instead, we had frequent discussions and exchanged ideas until the research topic started to take shape. I believe this is a valuable experience for becoming a scholar, as we need to come up with our own questions and effectively express our ideas to others in academia. This project also earned me the undergraduate scholarship funded by the National Science and Technology Council (NSTC) in Taiwan.

% After graduation, I became a master’s student and continued doing research in Prof.~Huang’s group and started collaborating with Dr.~Sungkit Yip, a distinguished research fellow at Academia Sinica, Taiwan. Prof.~Huang introduced me to Dr.~Yip during the 2024 summer internship at Academia Sinica. During the internship, we studied the deeper entanglement structure of the Kitaev ground state, especially how the topological order forms from the many-body entanglement perspective. I learned an important research skill from Dr.~Yip in this project: \textit{always start from a simple example}. We started from the two-spin model to a 1D spin chain, and gradually extended the theory to the complete 2D model. Each layer provides valuable insights into how spins get entangled and ultimately leads to the complete pattern. I also applied numerical skills acquired from undergraduate courses, such as symmetry-resolved exact diagonalization and the Lanczos method, to solve the Kitaev model, which enhanced my ability to analyse physics numerically. 

% Later, we switched our topic to the density matrix of Kitaev spin liquids because the density matrix is a powerful tool to study entanglement structure and may shed light on the quantum information description of a condensed matter system. It turned out that our work successfully revealed the organized structure of the density matrix of Kitaev spin liquids, allowing us to fully understand their entanglement spectrum. In this program, I mastered lattice gauge theory, higher-form symmetry, and quantum information-oriented tools used in many-body spin systems. This was also an important experience to practically apply my knowledge of quantum information and condensed matter acquired in university courses to real research. We wrote up the results into a paper, and it has now been published in Physical Review B. 

% At this stage, I am more eager to explore more possibilities and a broader range of condensed matter theory beyond my work. This desire makes UCSB my first choice due to the strong condensed matter theory group, where professors have diverse research interests. My research interests also align with professors at UCSB, including Leon Balents, Cenke Xu, and Chetan Nayak, ranging from frustrated magnets to quantum information. Besides, the Kavli Institute for Theoretical Physics (KITP) on campus is also one of the main reasons that attracts me. The international conferences and research opportunities, such as the KITP Graduate Fellowship Program, provided by this top-tier institute will be valuable opportunities for me to engage with leading physicists worldwide and broaden my perspective on emerging directions in condensed matter theory and beyond.  

% In summary, I believe my solid theoretical foundation and research experience qualify me to be a PhD student at UCSB. I also believe the infinite opportunities at UCSB will broaden my horizons and ultimately lead me to my goal of becoming an outstanding physicist. 